
\documentclass{article}
\usepackage{colortbl}
\usepackage{makecell}
\usepackage{multirow}
\usepackage{supertabular}

\begin{document}

\newcounter{utterance}

\centering \large Interaction Transcript for game `dond', experiment `coop\_en', episode 7 with qwen3.5{-}27b.
\vspace{24pt}

{ \footnotesize  \setcounter{utterance}{1}
\setlength{\tabcolsep}{0pt}
\begin{supertabular}{c@{$\;$}|p{.15\linewidth}@{}p{.15\linewidth}p{.15\linewidth}p{.15\linewidth}p{.15\linewidth}p{.15\linewidth}}
    \# & $\;$A & \multicolumn{4}{c}{Game Master} & $\;\:$B\\
    \hline

    \theutterance \stepcounter{utterance}  
    & & \multicolumn{4}{p{0.6\linewidth}}{
        \cellcolor[rgb]{0.9,0.9,0.9}{
            \makecell[{{p{\linewidth}}}]{
                \texttt{\tiny{[P1$\langle$GM]}}
                \texttt{You are playing a cooperative negotiation game in which you have to agree on how to divide a set of items among you and another player.} \\
\\ 
\texttt{Rules:} \\
\texttt{(a) You and the other player are given a set of items. Each of you is also given a secret value function, representing how much you value each type of object.} \\
\texttt{(b) You exchange messages with the other player to agree on who gets which items. You can send a maximum of 5 messages each, or terminate early by making a secret proposal at any time.} \\
\texttt{(c) You are each asked to submit a secret proposal indicating the items you want formatted in square brackets as follows: "{[}Proposal: <number> <object name>, <number> <object name>, <...>{]}"} \\
\texttt{(d) If your proposals are complementary, i.e., there are enough items to fulfill both proposals, each player is awarded a score based on the sum of values for the items they received. Otherwise, both of you get zero points.} \\
\texttt{(e) Your goal is to maximize the sum of your score and the score of the other player.} \\
\\ 
\texttt{Let us start.} \\
\\ 
\texttt{The set of available items is:} \\
\\ 
\texttt{1 paper, 3 pictures, 1 stone, 1 fork} \\
\\ 
\texttt{Your secret value function is:} \\
\\ 
\texttt{paper: 1, picture: 2, stone: 2, fork: 1} \\
\\ 
\texttt{IMPORTANT: Your messages, unless it is the secret proposal, are directly transmitted to the other player, so do not include any response to the rules or text announcing your message. To make a secret proposal, use the indicated format. Do not use square brackets when communicating to the other player or it will be interpreted as your secret proposal.} \\
            }
        }
    }
    & \\ \\

    \theutterance \stepcounter{utterance}  
    & \multicolumn{4}{p{0.6\linewidth}}{
        \cellcolor[rgb]{0.9,0.9,0.9}{
            \makecell[{{p{\linewidth}}}]{
                \texttt{\tiny{[P1$\rangle$GM]}}
                \texttt{Hi! Let's figure out the best way to split these items to maximize our combined score. I have my own values, but I'm sure you have yours too.} \\
\\ 
\texttt{Here is what I have:} \\
\texttt{{-} 1 paper} \\
\texttt{{-} 3 pictures} \\
\texttt{{-} 1 stone} \\
\texttt{{-} 1 fork} \\
\\ 
\texttt{My values are: paper (1), picture (2), stone (2), fork (1).} \\
\\ 
\texttt{Could you tell me which items you value the most? That way, we can give the high{-}value items to the person who wants them more, ensuring we get the highest total points possible. For example, if you value the stone more than I do, I'd be happy to let you have it.} \\
            }
        }
    }
    & & \\ \\

    \theutterance \stepcounter{utterance}  
    & & & \multicolumn{2}{p{0.3\linewidth}}{
        \cellcolor[rgb]{0.9,0.9,0.9}{
            \makecell[{{p{\linewidth}}}]{
                \texttt{\tiny{[GM$|$GM]}}
                \texttt{continue} \\
            }
        }
    }
    & & \\ \\

    \theutterance \stepcounter{utterance}  
    & & & \multicolumn{2}{p{0.3\linewidth}}{
        \cellcolor[rgb]{0.9,0.9,0.9}{
            \makecell[{{p{\linewidth}}}]{
                \texttt{\tiny{[GM$|$GM]}}
                \texttt{Hi! Let's figure out the best way to split these items to maximize our combined score. I have my own values, but I'm sure you have yours too.} \\
\\ 
\texttt{Here is what I have:} \\
\texttt{{-} 1 paper} \\
\texttt{{-} 3 pictures} \\
\texttt{{-} 1 stone} \\
\texttt{{-} 1 fork} \\
\\ 
\texttt{My values are: paper (1), picture (2), stone (2), fork (1).} \\
\\ 
\texttt{Could you tell me which items you value the most? That way, we can give the high{-}value items to the person who wants them more, ensuring we get the highest total points possible. For example, if you value the stone more than I do, I'd be happy to let you have it.} \\
            }
        }
    }
    & & \\ \\

    \theutterance \stepcounter{utterance}  
    & & \multicolumn{4}{p{0.6\linewidth}}{
        \cellcolor[rgb]{0.9,0.9,0.9}{
            \makecell[{{p{\linewidth}}}]{
                \texttt{\tiny{[GM$\rangle$P2]}}
                \texttt{You are playing a cooperative negotiation game in which you have to agree on how to divide a set of items among you and another player.} \\
\\ 
\texttt{Rules:} \\
\texttt{(a) You and the other player are given a set of items. Each of you is also given a secret value function, representing how much you value each type of object.} \\
\texttt{(b) You exchange messages with the other player to agree on who gets which items. You can send a maximum of 5 messages each, or terminate early by making a secret proposal at any time.} \\
\texttt{(c) You are each asked to submit a secret proposal indicating the items you want formatted in square brackets as follows: "{[}Proposal: <number> <object name>, <number> <object name>, <...>{]}"} \\
\texttt{(d) If your proposals are complementary, i.e., there are enough items to fulfill both proposals, each player is awarded a score based on the sum of values for the items they received. Otherwise, both of you get zero points.} \\
\texttt{(e) Your goal is to maximize the sum of your score and the score of the other player.} \\
\\ 
\texttt{Let us start.} \\
\\ 
\texttt{The set of available items is:} \\
\\ 
\texttt{1 paper, 3 pictures, 1 stone, 1 fork} \\
\\ 
\texttt{Your secret value function is:} \\
\\ 
\texttt{paper: 2, picture: 2, stone: 1, fork: 1} \\
\\ 
\texttt{IMPORTANT: Your messages, unless it is the secret proposal, are directly transmitted to the other player, so do not include any response to the rules or text announcing your message. To make a secret proposal, use the indicated format. Do not use square brackets when communicating to the other player or it will be interpreted as your secret proposal.} \\
\\ 
\\ 
\texttt{Hi! Let's figure out the best way to split these items to maximize our combined score. I have my own values, but I'm sure you have yours too.} \\
\\ 
\texttt{Here is what I have:} \\
\texttt{{-} 1 paper} \\
\texttt{{-} 3 pictures} \\
\texttt{{-} 1 stone} \\
\texttt{{-} 1 fork} \\
\\ 
\texttt{My values are: paper (1), picture (2), stone (2), fork (1).} \\
\\ 
\texttt{Could you tell me which items you value the most? That way, we can give the high{-}value items to the person who wants them more, ensuring we get the highest total points possible. For example, if you value the stone more than I do, I'd be happy to let you have it.} \\
            }
        }
    }
    & \\ \\

    \theutterance \stepcounter{utterance}  
    & & & \multicolumn{4}{p{0.6\linewidth}}{
        \cellcolor[rgb]{0.9,0.9,0.9}{
            \makecell[{{p{\linewidth}}}]{
                \texttt{\tiny{[GM$\langle$P2]}}
                \texttt{Hi! Thanks for sharing your values. It looks like we have a great opportunity to maximize our combined score by trading based on our differences.} \\
\\ 
\texttt{Here are my values:} \\
\texttt{{-} Paper: 2} \\
\texttt{{-} Picture: 2} \\
\texttt{{-} Stone: 1} \\
\texttt{{-} Fork: 1} \\
\\ 
\texttt{Comparing our lists:} \\
\texttt{{-} **Paper**: I value it at 2, you value it at 1. I should get the paper.} \\
\texttt{{-} **Pictures**: We both value them at 2. It doesn't matter who gets them for the total score, but since there are 3, we can split them or give them all to one person if it helps balance the deal.} \\
\texttt{{-} **Stone**: You value it at 2, I value it at 1. You should get the stone.} \\
\texttt{{-} **Fork**: We both value it at 1.} \\
\\ 
\texttt{To maximize the total points:} \\
\texttt{{-} I take the **paper** (gaining 2 points for me vs 1 for you).} \\
\texttt{{-} You take the **stone** (gaining 2 points for you vs 1 for me).} \\
\texttt{{-} For the **pictures** and **fork**, since our values are identical, we can just split them however we like. Let's say I take 2 pictures and you take 1 picture and the fork, or we split the pictures 2{-}1 and I take the fork. Actually, since the values are the same, the specific split of pictures and fork doesn't change the *total* sum, only who gets what.} \\
\\ 
\texttt{Let's propose a split:} \\
\texttt{{-} **Me**: 1 paper, 2 pictures. (My score: 2 + 2 + 2 = 6)} \\
\texttt{{-} **You**: 1 stone, 1 picture, 1 fork. (Your score: 2 + 2 + 1 = 5)} \\
\texttt{Total combined score: 11.} \\
\\ 
\texttt{Alternatively, if you prefer the fork, we can swap. Does this sound good to you? If so, I can make a formal proposal.} \\
\\ 
\texttt{{[}Proposal: 1 paper, 2 pictures{]}} \\
            }
        }
    }
     \\ \\

    \theutterance \stepcounter{utterance}  
    & & & \multicolumn{2}{p{0.3\linewidth}}{
        \cellcolor[rgb]{0.9,0.9,0.9}{
            \makecell[{{p{\linewidth}}}]{
                \texttt{\tiny{[GM$|$GM]}}
                \texttt{proposal} \\
            }
        }
    }
    & & \\ \\

    \theutterance \stepcounter{utterance}  
    & & & \multicolumn{2}{p{0.3\linewidth}}{
        \cellcolor[rgb]{0.9,0.9,0.9}{
            \makecell[{{p{\linewidth}}}]{
                \texttt{\tiny{[GM$|$GM]}}
                [1, 2, 0, 0]
            }
        }
    }
    & & \\ \\

    \theutterance \stepcounter{utterance}  
    & & \multicolumn{4}{p{0.6\linewidth}}{
        \cellcolor[rgb]{0.9,0.9,0.9}{
            \makecell[{{p{\linewidth}}}]{
                \texttt{\tiny{[P1$\langle$GM]}}
                \texttt{The other player has submitted a secret proposal. It is now your turn to submit a proposal of your own.} \\
            }
        }
    }
    & \\ \\

    \theutterance \stepcounter{utterance}  
    & \multicolumn{4}{p{0.6\linewidth}}{
        \cellcolor[rgb]{0.9,0.9,0.9}{
            \makecell[{{p{\linewidth}}}]{
                \texttt{\tiny{[P1$\rangle$GM]}}
                \texttt{{[}Proposal: 1 paper, 3 pictures, 1 stone, 1 fork{]}} \\
            }
        }
    }
    & & \\ \\

    \theutterance \stepcounter{utterance}  
    & & & \multicolumn{2}{p{0.3\linewidth}}{
        \cellcolor[rgb]{0.9,0.9,0.9}{
            \makecell[{{p{\linewidth}}}]{
                \texttt{\tiny{[GM$|$GM]}}
                \texttt{proposal} \\
            }
        }
    }
    & & \\ \\

    \theutterance \stepcounter{utterance}  
    & & & \multicolumn{2}{p{0.3\linewidth}}{
        \cellcolor[rgb]{0.9,0.9,0.9}{
            \makecell[{{p{\linewidth}}}]{
                \texttt{\tiny{[GM$|$GM]}}
                [1, 3, 1, 1]
            }
        }
    }
    & & \\ \\

    \theutterance \stepcounter{utterance}  
    & & & \multicolumn{2}{p{0.3\linewidth}}{
        \cellcolor[rgb]{0.9,0.9,0.9}{
            \makecell[{{p{\linewidth}}}]{
                \texttt{\tiny{[GM$|$GM]}}
                [[1, 3, 1, 1], [1, 2, 0, 0]]
            }
        }
    }
    & & \\ \\

\end{supertabular}
}

\end{document}
